\maketitle

\begin{abstract}
eQTL prediction (eQTLP) is the only variant task in DNALongBench
\citep{cheng2025dnalongbench}, a standard benchmark for long-context DNA models.
Ranking variants by distance to the transcription start site (TSS), with no
training, reaches a mean test AUROC of 0.751, beating the expert model (0.681)
and every foundation model (FM; 0.51--0.57). The test sets hold 11--34 positives
per tissue, and every test gene also appears in training. We rebuild the task
with distance-matched negatives and chromosome holdout, where hand-crafted local
sequence features set a floor of 0.599. Scoring each pair again with the
alternate allele replaced by the reference shows that no model uses the allele:
for a CNN trained end to end and for frozen HyenaDNA and Caduceus, the swap
changes AUROC by at most 0.0017, and all three fall below the floor. Fine-tuning
HyenaDNA end to end over 3 folds does not change this: accuracy improves to the
frozen probe's level while the allele gap stays at $|\Delta| \le 0.00016$. A nonlinear
probe on the allele difference alone recovers only a weak trace (0.52--0.55). A
single-nucleotide change moves the pooled FM embedding by about $10^{-4}$ of its
norm, and the released FM eQTL code slices every input to length zero. On 652
ClinVar pathogenic regulatory SNVs matched to benign SNVs, neither the embedding
shift nor the transferred probe separates the classes (AUROC 0.49--0.51). We
release the rebuilt task, the ablation harness, and the matched ClinVar set.
\end{abstract}

\section{Introduction}

Long-context DNA models \citep{nguyen2023hyenadna,schiff2024caduceus,brixi2026evo2}
are motivated by distal regulation, meaning variants that act on genes tens to
hundreds of kilobases away. They inherit their architectures from long-convolution
and state-space sequence models \citep{poli2023hyena,gu2023mamba}, and are
benchmarked against expert models such as Enformer \citep{avsec2021enformer},
Akita \citep{fudenberg2020akita} and Puffin \citep{dudnyk2024puffin}.
DNALongBench \citep{cheng2025dnalongbench} evaluates HyenaDNA
\citep{nguyen2023hyenadna} and Caduceus \citep{schiff2024caduceus} at up to
450\,kb. Of its five tasks, only eQTLP involves a variant. Its model interface
embeds the reference and alternate sequences, averages each, concatenates them,
and classifies. A model can therefore succeed without distinguishing the alleles
at all, if \emph{where} the variant sits predicts the label.

We show that this is what happens. The published split rewards distance and gene
identity. Once both are removed, three architectures score the same with and
without the alternate allele. We contribute: (1) an audit of the published
protocol; (2) \textbf{eQTLP-v2}, a distance-matched, chromosome-held-out version
with a published floor; (3) a \textbf{ref-copy ablation} that any variant
benchmark can report; (4) reproducibility findings on the released code; and (5)
a matched, evaluation-only ClinVar regulatory set.

\section{The published protocol rewards distance and memorisation}

Positives are SuSiE \citep{wang2020susie} fine-mapped GTEx \citep{gtex2020} eQTLs
in 9 tissues. Negatives are other variant--gene pairs within 450\,kb. Each input
spans the variant and the TSS and is right-padded with N to 450\,kb. The split is
a per-tissue random 8:1:1.

\begin{itemize}
\item \textbf{Distance beats every model.} Positives lie a median 8--33\,kb from
the TSS, negatives 88--145\,kb. Scoring by $-$distance ($|$variant $-$ TSS$|$
from the released coordinates, which is what the padded input encodes) gives a
test AUROC of 0.688--0.828 per tissue, with a mean of 0.751 (Table~\ref{tab:published}).
That beats the expert model in 7 of 9 tissues and every FM in every tissue. The
padding also hands the distance to the model, since the number of N tokens
encodes it.
\ifdefined\RECOMB
Scoring by the released \texttt{distance\_to\_tss} column instead --- which
disagrees with the coordinates by over 1\,kb for 16\% of negatives and no
positives --- gives 0.663: still above every FM, though not the expert.
\fi
\item \textbf{The test sets cannot rank models.} They hold 11--34 positives per
tissue, giving AUROC CI half-widths of $\pm$0.10--0.18 \citep{hanley1982auc}, so
no published FM result is distinguishable from 0.5. Nor is the split stratified:
positive rates are 9--21\% in train, 3--5\% in test.
\item \textbf{Train and test overlap.} Every test gene appears in training, and
82--92\% of test pairs have a training variant within 100\,kb.
\end{itemize}

\ifdefined\RECOMB
Table~\ref{tab:tissues} gives the per-tissue picture. The distance baseline is
not winning on average while losing in places: it beats every foundation model in
every one of the nine tissues, and the expert model in seven. The confidence
intervals are the second half of the problem. With 11--34 positives a test AUROC
carries a 95\% CI roughly 0.10--0.19 wide, so the published FM numbers
(0.47--0.60 across tissues) are not separable from each other, from 0.5, or from
the expert model. A benchmark cannot rank models at this sample size whatever the
models do.
% generated by analysis/make_appendix_tables.py -- do not edit by hand
\begin{table}[tb]
\centering
\caption{Published eQTLP test split, per tissue. \textbf{dist} is our zero-training baseline, scored from the released coordinates, with a 95\% bootstrap CI; \textbf{expert} (Enformer), \textbf{hyena} and \textbf{cad-ph} are the benchmark's published numbers. Distance wins in every tissue, the CIs are 0.10--0.19 wide, and every test gene appears in training. Full columns in Appendix~A.}
\label{tab:tissues}
\footnotesize
\setlength{\tabcolsep}{3pt}
\begin{tabular}{lcccccccc}
\toprule
Tissue & $n_+$ & $n_-$ & dist & 95\% CI & expert & hyena & cad-ph & gene in train \\
\midrule
Adipose\_Subcutaneous & 16 & 316 & 0.715 & [0.55, 0.87] & 0.736 & 0.513 & 0.541 & 100\% \\
Artery\_Tibial & 27 & 542 & 0.744 & [0.64, 0.84] & 0.741 & 0.479 & 0.547 & 100\% \\
Cells\_Cultured\_fibroblasts & 11 & 350 & 0.688 & [0.54, 0.82] & 0.639 & 0.584 & 0.597 & 100\% \\
Muscle\_Skeletal & 16 & 397 & 0.798 & [0.68, 0.90] & 0.621 & 0.487 & 0.538 & 100\% \\
Nerve\_Tibial & 34 & 609 & 0.818 & [0.76, 0.88] & 0.683 & 0.511 & 0.588 & 100\% \\
Skin\_Not\_Sun\_Exposed\_Suprapubic & 17 & 456 & 0.828 & [0.73, 0.91] & 0.710 & 0.471 & 0.586 & 100\% \\
Skin\_Sun\_Exposed\_Lower\_leg & 26 & 607 & 0.694 & [0.60, 0.78] & 0.700 & 0.544 & 0.574 & 100\% \\
Thyroid & 26 & 524 & 0.699 & [0.58, 0.81] & 0.612 & 0.529 & 0.527 & 100\% \\
Whole\_Blood & 17 & 305 & 0.776 & [0.66, 0.88] & 0.689 & 0.512 & 0.594 & 100\% \\
\bottomrule
\end{tabular}
\end{table}

\fi

\begin{table}[tb]
\centering
\caption{Published test AUROC (Table 7 of the benchmark paper), mean over 9
tissues, next to our zero-training baseline. Per-tissue values with CIs are in
Appendix~A.}
\label{tab:published}
\small
\begin{tabular}{cccccc}
\toprule
Distance only & Expert (Enformer) & CNN & HyenaDNA & Caduceus-Ph & Caduceus-PS \\
\midrule
\textbf{0.751} & 0.681 & 0.528 & 0.514 & 0.566 & 0.538 \\
\bottomrule
\end{tabular}
\end{table}

\section{eQTLP-v2}

We pool the tissues into 25,600 unique pairs and drop 111 with conflicting
labels. Each positive is matched to one negative on log distance to the TSS, on
the same side of it, which gives 2,905 pairs and drops the distance-only AUROC to
0.495. Folds are 5-fold with whole chromosomes held out, giving 570--665 test
positives per fold. We also drop the published label-derived mask, which covers
the positions of the positive eQTLs.

Matching on distance removes real regulatory signal along with the shortcut,
since enhancer--promoter distance genuinely predicts effect. We match anyway,
because the task's stated aim is variant-effect prediction from sequence, and
because the published input encodes distance in its padding, so a model can
exploit it without reading any sequence. Unmatched results are in Appendix~B.

On this configuration, gradient-boosted trees on local composition (GC, CpG,
repeats and dinucleotides in windows up to 5\,kb) score $\mathbf{0.599 \pm
0.031}$. Adding allele features changes this by at most 0.01. On the random
split, a gene's training-set positive rate alone scores 0.70. That is the
memorisation the published split allows.

\ifdefined\RECOMB
Table~\ref{tab:shortcuts} separates the two shortcuts by turning each off
independently. Distance scores 0.790 whatever the split, because it does not care
which genes are held out --- and 0.500 once negatives are matched on it. The
gene prior behaves in the opposite way: 0.703 when genes recur across the split
and 0.387 under chromosome holdout, which is below chance because a gene's
training rate actively misleads on unseen chromosomes. The two are independent,
which is why removing only one leaves the task solvable without reading sequence.

The table also quantifies a third problem. The published label mask marks the
positions of the positive eQTLs, so simply counting masked bases scores 0.780 on
the unmatched task, above every published model. It falls to 0.436 once the
positives are distance-matched, confirming it is a label leak rather than signal.
We drop the mask in eQTLP-v2 for this reason.
% generated by analysis/make_appendix_tables.py -- do not edit by hand
\begin{table}[tb]
\centering
\caption{What each split lets a simple feature exploit, on eQTLP-v2. \textbf{all} keeps the original distance imbalance, \textbf{matched} removes it; \textbf{random} lets genes recur across the split, \textbf{chrom} does not. Distance falls 0.79$\to$0.50 under matching, and the gene prior falls 0.70$\to$0.44 under chromosome holdout: those two columns are the shortcuts. The label-mask count collapses from 0.78 to 0.44, confirming it is a label leak. Composition on matched/chrom is the floor; the 0.599 in the text is its mean over 7 matchings $\times$ 10 partitions, while this cell is the single seed-0 matching.}
\label{tab:shortcuts}
\small
\begin{tabular}{lcccc}
\toprule
Feature & all / random & all / chrom & matched / random & matched / chrom \\
\midrule
distance to TSS & 0.790 & 0.790 & 0.500 & 0.500 \\
gene's train positive rate & 0.703 & 0.387 & 0.706 & 0.439 \\
distance + label-mask count & 0.810 & 0.780 & 0.529 & 0.436 \\
local composition & 0.744 & 0.639 & 0.665 & 0.588 \\
distance + composition & 0.833 & 0.790 & 0.670 & 0.582 \\
distance + composition + allele & 0.834 & 0.790 & 0.670 & 0.590 \\
\bottomrule
\end{tabular}
\end{table}

\fi

\section{No model uses the allele}

\paragraph{Protocol.}
\ifdefined\RECOMB
The experiment is a controlled substitution. We score every variant twice with
the same already-fitted model: once with its real alternate allele, and once with
that allele overwritten by the reference base, so the two inputs differ at
exactly one position and in nothing else. A model that reads the variant must
give the two different scores. A model that is really scoring the variant's
location will give the same score both times.

Concretely, each model converts a sequence into a vector: it produces one
internal state per base, and those states are averaged across all positions into
a single summary vector $e$. The benchmark's own interface then concatenates the
reference and alternate summaries and classifies the pair. We report three
quantities from it:
\begin{itemize}
\item \refalt{} --- the test AUROC on the real pair $[e_{\mathrm{ref}},
e_{\mathrm{alt}}]$;
\item \refcopy{} --- \emph{the same fitted classifier}, nothing refitted, applied
to $[e_{\mathrm{ref}}, e_{\mathrm{ref}}]$, so the alternate allele is gone;
\item \diffonly{} --- a classifier trained only on the difference
$e_{\mathrm{alt}} - e_{\mathrm{ref}}$, which asks whether the allele's effect is
recoverable at all.
\end{itemize}

The models are the benchmark's \textbf{CNN} encoder, trained end to end with its
collapsing logit ReLU repaired (\S\ref{sec:code}); and
\textbf{HyenaDNA-medium-450k} \citep{nguyen2023hyenadna} (which reads sequence in
one direction only, 6.55\,M parameters) and \textbf{Caduceus-Ph-131k}
\citep{schiff2024caduceus} (which reads both directions, 7.73\,M). The latter two
are used \emph{frozen}: their pretrained weights are never updated, and we fit
only a logistic regression on top of the summary vectors they produce
\citep{pedregosa2011sklearn} --- a \emph{linear probe}, in the
representation-learning sense of that word rather than the molecular one. This is
standard practice for asking what
information a pretrained representation already contains, and we report below
what happens when we instead fine-tune the whole network. Inputs are
unpadded, and for each model we try both the published sequence orientation and
its reverse complement.
\else
Each model embeds the reference and alternate sequences.
The last-layer states are mean-pooled, and $[e_{\mathrm{ref}}, e_{\mathrm{alt}}]$
is classified, as in the benchmark's Methods. We report three quantities:
\refalt{}, the test AUROC; \refcopy{}, the same fitted classifier scoring
$[e_{\mathrm{ref}}, e_{\mathrm{ref}}]$; and \diffonly{}, a classifier trained
only on $e_{\mathrm{alt}} - e_{\mathrm{ref}}$. The models are the benchmark's
\textbf{CNN} encoder, trained end to end with its collapsing logit ReLU repaired
(\S\ref{sec:code}); and \textbf{HyenaDNA-medium-450k} \citep{nguyen2023hyenadna}
(causal, 6.55\,M parameters) and \textbf{Caduceus-Ph-131k}
\citep{schiff2024caduceus} (bidirectional, 7.73\,M), frozen, with a
logistic-regression probe \citep{pedregosa2011sklearn}. Inputs are unpadded. For
each FM we embed both the published orientation and its reverse complement.
\fi


\begin{itemize}
\item \textbf{The allele changes nothing} (Figure~\ref{fig:one}a). Replacing it
with the reference changes AUROC by at most 0.0017 for the trained CNN (15 runs)
and 0.0006 for the FM probes, across probe strengths ($C = 0.001$--10), heads
(linear, MLP), orientations, and 13 fold partitions (chromosome, gene and
random). The paired 95\% CI is $[-0.0002, 0.0002]$.
\item \textbf{Only a weak trace of the allele is there to find, and pooling is
not what hides it.} A linear \diffonly{} probe scores at chance (0.49--0.51). An
MLP on the difference beats a label-permutation null (Caduceus 0.54--0.55,
HyenaDNA 0.52--0.53) but stays below the floor, and an MLP head on
$[e_{\mathrm{ref}}, e_{\mathrm{alt}}]$, like the released Caduceus head, still
has $\Delta = 0.0000$.

Averaging over the whole sequence discards position, so we next ask whether the
allele survives somewhere that average cannot reach. We repeat the search at
every layer and over four \emph{summary windows} --- the set of positions
averaged together before classifying: the whole sequence, $\pm$1\,kb around the
variant, $\pm$64\,bp, and the single variant position itself. Each is scored
against a control in which the variant is replaced by a base that did
\textbf{not} occur, which absorbs anything attributable to local sequence context
rather than to the allele. For HyenaDNA the control matches the signal
everywhere (mean $-0.006$ over 36 cells), so what looked like allele information
was context. For Caduceus one window survives its control --- the final layer at
$\pm$1\,kb, scoring 0.575--0.580 against 0.545--0.553, on all three fold schemes.
So a variant-centred summary does recover allele signal that whole-sequence
averaging loses. It is still below the floor, and the swap still changes nothing
even there (max $|\Delta|$ 0.017 over 104 cells, always favouring \refcopy{}).
\item \textbf{All three models score below the hand-feature floor}
(Figure~\ref{fig:one}a), in 10 of 10 random chromosome partitions. Only on the
random split, which allows locus memorisation, do the FMs pass it (0.69--0.70 vs
0.67).
\end{itemize}

\ifdefined\RECOMB
Figure~\ref{fig:two} shows that search in full. Each cell is one layer read out
one way, scored by how much a nonlinear probe on the allele difference beats the
same probe on a control in which the variant is replaced by a base that never
occurred. The control is what makes the comparison informative: it absorbs
everything the probe could learn from local sequence context, so what remains is
attributable to the allele itself. Three things are visible. HyenaDNA has no cell
that clears its control, which is why we describe its trace as context rather
than allele. Caduceus is warm in a band --- the $\pm$1\,kb and $\pm$64\,bp
readouts, across most of the depth --- but almost all of it is within noise of
the control. And exactly one cell stands clear, the final layer at $\pm$1\,kb,
which replicates on all three fold schemes. That cell is the strongest
allele-specific signal anywhere in either model, and at 0.575--0.580 it is still
below the 0.599 hand-feature floor.

\fi

\begin{table}[tb]
\centering
\caption{eQTLP-v2, matched pairs, chromosome holdout. Mean per-fold test AUROC
$\pm$ SD.}
\label{tab:v2}
\small
\begin{tabular}{lcccc}
\toprule
Model & \refalt{} & \refcopy{} & $\Delta$ & \diffonly{} \\
\midrule
CNN (trained, 3 seeds $\times$ 5 folds) & 0.517 $\pm$ .040 & 0.517 & $-$0.0002 & --- \\
HyenaDNA (frozen)    & 0.556 $\pm$ .031 & 0.556 $\pm$ .030 & $-$0.0001 & 0.495 $\pm$ .021 \\
Caduceus-Ph (frozen) & 0.564 $\pm$ .021 & 0.564 $\pm$ .021 & 0.0000 & 0.507 $\pm$ .016 \\
Local-composition floor & 0.599 $\pm$ .031 & & & \\
\bottomrule
\end{tabular}
\end{table}

\paragraph{The SNP barely registers} (Figure~\ref{fig:one}b). The median relative
shift $\|e_{\mathrm{alt}} - e_{\mathrm{ref}}\| / \|e_{\mathrm{ref}}\|$ is
1.5--2.0 $\times 10^{-4}$ for HyenaDNA and $0.7 \times 10^{-4}$ for Caduceus.
eQTLP-v2 is 92.3\% single-nucleotide substitutions and 7.7\% short indels or
MNVs; restricting to the substitutions alone leaves both medians unchanged
($1.52 \times 10^{-4}$ and $0.674 \times 10^{-4}$), so this is a per-variant
effect and not an artifact of the indels. The
CNN's max-pooled embedding changes in only 0.4--1.7\% of pairs. Two explanations
for this fail.

\ifdefined\RECOMB
\paragraph{Orientation is not the cause.} HyenaDNA is causal, and the published
orientation places the variant roughly 500\,bp from the end of the input, so in
71--86\% of pairs only a few hundred positions can attend to it. This is a real
defect of the protocol, and it is the obvious suspect. It is not the explanation.
Embedding the reverse complement instead, so that the variant comes first and
every subsequent position can see it, moves AUROC by less than 0.01 and leaves
the \refcopy{} gap in the fourth decimal. Caduceus is bidirectional by
construction and behaves identically to HyenaDNA throughout. A null that survives
both reading orders, in both a causal and a bidirectional architecture, is not an
artifact of where the variant sits in the window.

\paragraph{Numerical precision is not the cause.} A relative shift of $10^{-4}$
in bf16 invites the objection that the allele is being rounded away. We
re-embedded 400 pairs in full fp32: the shifts reproduce, the difference vectors
agree with their bf16 counterparts at cosine 0.89--0.95, and \diffonly{} stays at
chance. The signal that a probe cannot find in bf16 is equally absent at full
precision, so the model is not computing an allele effect that the numerics
discard.

\paragraph{What remains is pooling.} The substitution does move the
representation, and locally it moves it a great deal: the variant-position hidden
state changes by 37--112\%, the $\pm$64\,bp mean by 0.7--6.1\%, and the
whole-sequence mean by only 0.006--0.09\%. Taking medians across layers, what
reaches the classifier is about $1{,}000\times$ smaller than the change at the
variant itself. Averaging across tens of kilobases is simply the wrong way to
summarise a one-base question, and it is the summary the benchmark specifies.
This is why we recommend summarising a window around the variant in
\S\ref{sec:rec} --- but note the limit of that recommendation:
Figure~\ref{fig:two} shows that even the best local window recovers only a signal
that remains below the hand-feature floor.
\else
\textbf{Orientation:} HyenaDNA is causal, and in the published
orientation most variants sit in the last 500\,bp; but reversing the input leaves
AUROC unchanged, and bidirectional Caduceus behaves identically.
\textbf{Precision:} re-embedding 400 pairs in fp32 gives the same shift
(bf16--fp32 difference-vector cosine 0.89--0.95), and \diffonly{} stays at
chance. What remains is pooling. The SNP moves the variant-position hidden state
by 37--112\%, the $\pm$64\,bp mean by 0.7--6.1\%, and the whole-sequence mean by
0.006--0.09\%: mean-pooling over tens of kilobases dilutes the substitution
${\sim}1{,}000\times$.
\fi

\begin{figure}[!htb]
\centering
\ifdefined\RECOMB
  \includegraphics[width=\linewidth]{figure1_wide.pdf}
\else
  \includegraphics[width=\linewidth]{figure1.pdf}
\fi
\caption{\textbf{(a)} eQTLP-v2, matched pairs, chromosome holdout: per-fold test
AUROC for each model, scored normally (\refalt{}, filled) and with the reference
copied over the alternate allele (\refcopy{}, open). Each fold's pair is joined
by a line, and every line is vertical: the swap moves no score. Points are 5
folds ($\times$ 3 seeds for the CNN); the bar is the mean. The dashed line is the
floor on these same folds (0.593); \S3 quotes $0.599 \pm 0.031$, its mean over 7
matchings $\times$ 10 partitions. No model reaches it. \textbf{(b)} Relative
shift
$\|e_{\mathrm{alt}} - e_{\mathrm{ref}}\| / \|e_{\mathrm{ref}}\|$ under four
summary windows, from the variant position out to a whole-sequence average; markers are the
median over layers, bands the min--max.}
\label{fig:one}
\end{figure}

\ifdefined\RECOMB
\begin{figure}[!htb]
\centering
\includegraphics[width=\linewidth]{figure2_wide.pdf}
\caption{Allele signal by layer and summary window, for both frozen models. Colour is
the gain of a nonlinear probe on $e_{\mathrm{alt}} - e_{\mathrm{ref}}$ over the
same probe on a random-allele control, so 0 means the probe was reading local
context rather than the variant. \textbf{F} is the pooled final-layer state.
Warm cells concentrate in the local windows and only one (Caduceus, final layer,
$\pm$1\,kb) clears its control by a margin; it is still below the floor, and the
\refcopy{} swap changes nothing there.}
\label{fig:two}
\end{figure}
\fi

\paragraph{Fine-tuning does not change it.} The probes above are frozen, whereas
the published numbers come from fine-tuning, so we fine-tuned
HyenaDNA-medium-450k end to end on eQTLP-v2 --- backbone and head, through the
benchmark's own mean-pooled readout --- for 2 epochs on each of 3 chromosome
folds (33\,GPU-hours; inputs capped at 131\,kb, which 88\% of pairs fit
entirely).

\ifdefined\RECOMB
Training works. Loss falls from 0.730 to 0.698, and \texttt{refalt} rises across
epochs on 2 of the 3 folds (fold 0: 0.543$\to$0.563; fold 1: 0.565$\to$0.572;
fold 2 regresses, 0.566$\to$0.529). The final-epoch mean is $0.555 \pm 0.023$,
statistically indistinguishable from the frozen probe's $0.556 \pm 0.031$ and
still below the 0.599 floor: updating every weight in the backbone buys nothing
over a logistic regression on its frozen features.

The allele gap does not move. Across all six epochs the largest $|\Delta|$ is
\textbf{0.00016}, with a mean of $+0.00005$ --- the same fourth-decimal nothing
the frozen probes gave. The model measurably improves at the task while learning
nothing about which allele is present. This is the strongest form of the result
we can offer: the null is not an artifact of freezing the backbone.

% generated by analysis/make_appendix_tables.py -- do not edit by hand
\begin{table}[tb]
\centering
\caption{End-to-end fine-tuning of HyenaDNA-medium-450k on eQTLP-v2, through the benchmark's own readout. Training works --- loss falls and \texttt{refalt} rises on 2 of 3 folds --- and lands at the frozen probe's accuracy, but the allele gap never leaves the fourth decimal.}
\label{tab:ft}
\small
\begin{tabular}{lccccccc}
\toprule
Fold & Epoch & \texttt{refalt} & \texttt{ref\_copy} & $\Delta$ & train loss & $n$ & h \\
\midrule
0 & 1 & 0.5430 & 0.5429 & +0.00008 & 0.7295 & 1212 & 4.6 \\
0 & 2 & 0.5632 & 0.5633 & -0.00004 & 0.7064 & 1212 & 6.5 \\
1 & 1 & 0.5651 & 0.5652 & -0.00001 & 0.7126 & 953 & 4.6 \\
1 & 2 & 0.5722 & 0.5721 & +0.00016 & 0.7020 & 953 & 5.4 \\
2 & 1 & 0.5663 & 0.5663 & +0.00005 & 0.7126 & 1083 & 5.4 \\
2 & 2 & 0.5293 & 0.5293 & +0.00002 & 0.6983 & 1083 & 5.5 \\
\bottomrule
\end{tabular}
\end{table}

\else
Training works --- loss falls 0.730$\to$0.698 and \refalt{} rises on 2 of 3
folds --- and lands at $0.555 \pm 0.023$, indistinguishable from the frozen
probe's $0.556 \pm 0.031$ and still below the floor. The allele gap does not
move: across all six epochs the largest $|\Delta|$ is \textbf{0.00016}. The model
improves at the task while learning nothing about which allele is present, so the
null is not an artifact of freezing the backbone.
\fi

\paragraph{Scope and limitations.} The fine-tuning run is one model, one seed, 2
epochs per fold at 131\,kb rather than the published 450\,kb, so it bounds what
this much training achieves, not what any training could. We evaluate two models
and both are small (6.55\,M and 7.73\,M parameters); larger long-context models
such as Evo 2 \citep{brixi2026evo2} are beyond our compute. Our claim is
therefore about what these architectures and this readout do on this task, not
about model scale.

\section{The released code does not run as shipped}
\label{sec:code}

\ifdefined\RECOMB
We report defects found by reading the released code and running it on real
records. We do not claim to have re-run the published experiments and obtained
different numbers, which would be a different kind of result: with these defects
the released code cannot produce the published FM numbers at all, so those
numbers must come from a version that is not the one released. We raise this as
a reproducibility problem for the artifact, not as a conclusion about what the
original authors ran.
\else
We report defects found by reading and running the released code; we do not claim
to have re-run the published experiments, since with these defects the released
code cannot produce the published FM numbers at all.
\fi

\begin{itemize}
\item \textbf{The FM eQTL loaders produce empty inputs.} Both the HyenaDNA and
Caduceus loaders pad each input to 450,000 positions, then slice it with
\texttt{[2028500:-2028500]}, leaving an empty array.
\item \textbf{The token ids are wrong.} Without the slice, \texttt{argmax} over
the one-hot gives ids 0--3, which are special tokens in the pretrained
vocabularies.
\item \textbf{The configs don't match the checkpoints.} They specify 2 layers,
while the checkpoints have 8 and 16, and the pretrained-weights path is unset.
\item \textbf{The CNN's training often collapses.} Its eQTLP head ends in a ReLU
on the logits, so both logits can clamp to zero: the prediction is constant and
every weight gets zero gradient. On the \emph{published} split, 2 of 10
initialisations are dead before the first step, and after one epoch 5--6 of 10
seeds output a constant score (AUROC exactly 0.5). The script sets no seed and
evaluates the final-epoch model, not the best. The CNN is fed only the reference.
\end{itemize}

\section{Clinical evaluation set}

\ifdefined\RECOMB
ClinVar \citep{landrum2025clinvar} pathogenic or likely-pathogenic regulatory
SNVs with at least 1 star, annotated against GENCODE \citep{gencode2025}, must be
non-coding, over 50\,bp from a splice site, off every CDS, and outside small
structural RNAs; 652 map to a gene TSS within 450\,kb. Each is matched to a
benign SNV passing the same filters, on distance to the TSS and, where possible,
the same gene and location class (555 of 652 pairs). No model is fitted on
clinical labels, and the embedding shift, the transferred eQTLP-v2 probe and that
probe's allele-only component are all at chance (AUROC 0.489--0.507 overall,
0.49--0.52 on the 285 pairs with a 2-star-or-better pathogenic). The only stratum
above chance is 5$'$UTR/promoter (0.53--0.57), where matching left pathogenic
variants closer to the TSS (median 170 vs 450\,bp) and distance alone scores
0.61. The same null therefore appears under pathogenicity labels; the two targets
differ, so this corroborates the eQTL result rather than generalising it. The set
is evaluation-only because too few regulatory positives survive the filters to
train on (409 at 2 stars).

The set is small and concentrated: a handful of well-studied genes contribute
many pairs each. Appendix~F repeats the evaluation after dropping the three most
frequent genes, after dropping the ten most frequent, and after keeping one pair
per gene. Every score stays at chance, so the null is not an artifact of which
genes ClinVar happens to have annotated densely.
\else
ClinVar \citep{landrum2025clinvar} pathogenic or likely-pathogenic regulatory
SNVs with at least 1 star are filtered to non-coding variants over 50\,bp from a
splice site, off every CDS and outside small structural RNAs; 652 map to a gene
TSS within 450\,kb, each matched to a benign SNV passing the same filters on
distance to the TSS and, where possible, gene and location class. No model is
fitted on clinical labels, and the embedding shift, the transferred eQTLP-v2
probe and its allele-only component are all at chance (0.489--0.507). The only
stratum above chance is 5$'$UTR/promoter (0.53--0.57), where matching left
pathogenic variants closer to the TSS and distance alone scores 0.61. The same
null therefore appears under pathogenicity labels; the two targets differ, so
this corroborates the eQTL result rather than generalising it.
\fi

\section{Relation to prior work}

Two long-range DNA benchmarks are adjacent to ours: the Genomics Long-Range
Benchmark \citep{kao2024lrb}, with an OMIM pathogenic non-coding variant task,
and BEND \citep{marin2024bend}, on reference-sequence element annotation. Neither
reports an allele ablation, the diagnostic we argue for: without it, a variant
task can be passed by locating the variant rather than reading it. Our result
also extends to a benchmark setting the finding that sequence-to-expression
models predict individual expression variation poorly
\citep{sasse2023benchmarking,huang2023personal}. The failure mode is shortcut
learning \citep{geirhos2020shortcut}: the protocol, not the architecture, decides
what a model must learn.

\ifdefined\RECOMB
\paragraph{Distance is a known baseline, not a surprise.} That genomic distance
predicts enhancer--gene links is long established: the activity-by-contact model
combines enhancer activity with contact frequency, which falls off steeply with
distance, and outperforms many sequence-based predictors of which enhancer
regulates which gene \citep{fulco2019abc}. The finding here is therefore not that
distance is informative --- it is that a benchmark advertised as measuring
variant-effect prediction from sequence can be won by distance alone, and that
none of the evaluated models beats it. A field that already knows distance is
strong should be reporting it as a baseline.

\paragraph{Perturbation data avoids the fine-mapping confound.} Our labels are
statistically fine-mapped eQTLs, which inherit the confounding structure of
linkage disequilibrium: the credible set marks variants correlated with an
expression effect, not necessarily the causal one. CRISPR interference screens
that perturb candidate elements directly and read out expression
\citep{gasperini2019,schraivogel2020} give element--gene links that do not
depend on fine-mapping assumptions, and they are a better foundation for a
variant benchmark than association statistics are --- though they test elements
rather than single bases, so they do not by themselves solve the allele problem
we identify.

\paragraph{Model scale.} We evaluate two small models. Larger genomic foundation
models exist --- Nucleotide Transformer \citep{dallatorre2025nt} at up to 2.5\,B
parameters, Evo 2 \citep{brixi2026evo2} at 7\,B and 40\,B --- and we have not
tested whether scale changes the picture. Two considerations cut against
assuming it does. The pooled readout that dilutes the substitution
${\sim}1{,}000\times$ is a property of the interface, not the backbone, so a
larger backbone read out the same way inherits the same dilution. And
pretraining on a single reference genome provides no objective that rewards
distinguishing alleles at a locus, whatever the parameter count. Neither
argument is a substitute for running the experiment.
\fi

\ifdefined\RECOMB
\section{Discussion}

\paragraph{What this does and does not show.} It shows that DNALongBench's eQTL
task, as published, can be passed without reading the alternate allele, and that
three architectures do exactly that. It does not show that long-context DNA
models are incapable of variant-effect prediction. The distinction matters
because the two conclusions call for different responses: the first is a
protocol that needs fixing, the second would be an architectural verdict we have
no basis to issue. Our evidence is two frozen models of 6.55\,M and 7.73\,M
parameters plus one CNN trained end to end, on one task, under the benchmark's
own readout. A larger model, a variant-aware readout, or a training signal that
rewards allele discrimination could all change the result --- and the fact that
none of them has been tested is a gap in the field, not a settled question.

\paragraph{Why a protocol can hide this.} The failure is legible only because we
looked for it. A benchmark reports one number per model, and 0.566 for
Caduceus-Ph is unremarkable: neither so high as to invite an audit nor so low as
to suggest the task is broken. Nothing in that number reveals that a
zero-training distance baseline scores 0.751 on the same split, that the test set
holds 27 positives, or that replacing the alternate allele with the reference
leaves the score unchanged. Each of those is a cheap check. None is standard.

\paragraph{The ablation generalises past this task.} The \refcopy{} construction
needs no retraining and no new data: take the fitted model, feed it an input in
which the quantity under test has been neutralised, and score again. It applies
wherever a benchmark claims a model responds to a specific perturbation --- a
variant, an edit, a knockout, a designed mutation. We think a variant benchmark
that does not report it has not established that its models read variants, and
that the burden should sit with the benchmark rather than with whoever audits it
afterwards.

\paragraph{On the value of floors.} The hand-feature floor did more diagnostic
work here than any model comparison. A floor of 0.599 from GC content, CpG
density, repeats and dinucleotide frequencies is not a strong result; what makes
it useful is that three deep models fail to reach it. Model-versus-model tables
cannot produce that observation, because every entry can be mediocre together.
Publishing a floor costs a few CPU-hours and converts an unremarkable
leaderboard into a falsifiable claim.
\fi

\section{Recommendations}
\label{sec:rec}

Variant benchmarks for long-context models should match negatives on distance to
the TSS and, where possible, gene; hold out chromosomes; size test sets for
$\pm$0.05 CIs; publish a hand-feature floor; report a \textbf{ref-copy ablation}
beside every variant score; and read the variant out \emph{locally} rather than
as a whole-input mean.

On DNALongBench's eQTL task, current models learn where a variant is, not what it
is.

\bibliographystyle{plainnat}
\bibliography{references}

\ifdefined\RECOMB
\section*{Appendix}
\noindent Appendices A--F are submitted as Supplementary Material
(\texttt{appendix.pdf}).
\else
\appendix
\section*{Appendix}
\noindent
\textbf{A} per-tissue Table~\ref{tab:published} with CIs, test-set sizes and
overlap statistics; \textbf{B} eQTLP-v2 construction and the full baseline grid
(random, gene and chromosome holdout; all vs matched); \textbf{C} per-fold
ablation results, results by orientation, and the fp32 check; \textbf{D} the CNN
collapse analysis; \textbf{E} the released-code defects with file and line
references; \textbf{F} clinical-set construction and results by stratum.
\fi
